\documentclass[11pt]{article}

\usepackage[margin=1in]{geometry}
\usepackage{amsmath,amssymb,amsthm,mathtools}
\usepackage{graphicx}
\usepackage[hidelinks]{hyperref}
\usepackage{microtype}

\newtheorem{theorem}{Theorem}
\newtheorem{lemma}[theorem]{Lemma}
\newtheorem{remark}[theorem]{Remark}
\newcommand{\T}{\mathbb T}
\newcommand{\R}{\mathbb R}
\newcommand{\Z}{\mathbb Z}
\newcommand{\cts}{\mathrm{cts}}
\newcommand{\Hom}{\operatorname{Hom}}

\title{Globally Continuous Cohomology of Tori\\
\large A full answer to Question 10.1 of arXiv:1305.7269}
\author{Candidate solution packet for human review}
\date{13 August 2026}

\begin{document}
\maketitle

\begin{abstract}
Austin asked whether the globally continuous bar-cochain group
$H^3_{\cts}(\T,\T)$ is nonzero.  We prove that it is zero.  More generally,
for the (necessarily trivial) continuous action of $\T^d$ on $\T^m$, we
compute the whole theory:
\[
H^n_{\cts}(\T^d,\T^m)\cong
\begin{cases}
\T^m,&n=0,\\
\operatorname{Mat}_{m\times d}(\Z),&n=1,\\
0,&n\ge2.
\end{cases}
\]
The proof avoids the unavailable long exact coefficient sequence.  A
normalized cochain in degree at least two factors through a simply connected
smash product, hence lifts to a real cochain.  For a cocycle the lift is a
real cocycle, and Haar averaging supplies its primitive.
\end{abstract}

\section{The published question}

Let $\T=\R/\Z$, written additively.  In Section 10.1 of
\cite{AustinPDE}, Austin considers the cohomology $H^*_{\cts}$ of the usual
inhomogeneous bar complex, restricted to globally continuous cochains.  The
paper observes that this theory need not provide long exact sequences for
general short exact sequences of topological modules, and asks:

\begin{quote}
\textbf{Question 10.1.} Is $H^3_{\cts}(\T,\T)$ non-zero?
\end{quote}

The question appears on physical page 91 of the source PDF:

\begin{center}
\includegraphics[width=0.94\textwidth]{figures/open_problem_crop.png}
\end{center}

Austin--Moore define exactly this theory as continuous maps in the ordinary
bar resolution \cite[Section 2]{AustinMoore}.  For a topological group $G$
acting trivially on an abelian topological group $A$, put
$C^n_{\cts}(G,A)=C(G^n,A)$ and
\begin{align*}
(dc)(g_1,\ldots,g_{n+1})
={}&c(g_2,\ldots,g_{n+1})\\
&+\sum_{i=1}^{n}(-1)^i
c(g_1,\ldots,g_i g_{i+1},\ldots,g_{n+1})\\
&+(-1)^{n+1}c(g_1,\ldots,g_n).
\end{align*}
Then $H^n_{\cts}(G,A)=\ker d/\operatorname{im}d$ in this complex.

\section{Result and proof intuition}

\begin{theorem}[Full torus computation]\label{thm:main}
Let $d,m\ge1$.  For the trivial action of $\T^d$ on $\T^m$,
\[
H^n_{\cts}(\T^d,\T^m)\cong
\begin{cases}
\T^m,&n=0,\\
\Hom_{\cts}(\T^d,\T^m)
 \cong\operatorname{Mat}_{m\times d}(\Z),&n=1,\\
0,&n\ge2.
\end{cases}
\]
In particular,
\[
\boxed{H^3_{\cts}(\T,\T)=0,}
\]
so the answer to Question 10.1 is no.
\end{theorem}

\paragraph{Idea of the proof.}
Normalize the bar complex.  A normalized $n$-cochain vanishes whenever one
input is the identity.  For $n\ge2$ it therefore factors through the smash
product $(\T^d)^{\wedge n}$.  This smash product is simply connected, so the
cochain has a continuous logarithm $C:(\T^d)^n\to\R^m$.  If the original
cochain is a cocycle, $dC$ is lattice-valued.  Connectedness makes it
constant, and normalization makes that constant zero.  Thus $C$ is an actual
real cocycle.  Finally, averaging its last variable against Haar measure gives
an explicit real primitive.  Reducing that primitive modulo $\Z^m$ completes
the argument.

\section{Three elementary lemmas}

\begin{lemma}[Continuous normalization]\label{lem:normalization}
The subcomplex of normalized continuous cochains
\[
N^n_{\cts}(G,A):=
\{c\in C^n_{\cts}(G,A):c(g_1,\ldots,g_n)=0
\text{ if some }g_i=e\}
\]
has the same cohomology as $C^*_{\cts}(G,A)$.
\end{lemma}

\begin{proof}
This is the usual normalization theorem for the bar complex.  Algebraically,
the bar cochains form a cosimplicial abelian group; the normalized Moore
complex is a deformation retract, while the complementary degenerate complex
is contracted by the codegeneracies.  In inhomogeneous coordinates those
codegeneracies insert the identity in one coordinate.  The projection and
contracting homotopy are finite sums of coface and codegeneracy operators.
All of these are compositions with continuous group operations or identity
insertions, so they preserve continuity.  Hence the standard algebraic
normalization proof restricts verbatim to globally continuous cochains.
\end{proof}

\begin{lemma}[Normalized torus cochains lift]\label{lem:lift}
Let $n\ge2$.  Every normalized continuous map
$c:(\T^d)^n\to\T^m$ has a normalized continuous lift
$C:(\T^d)^n\to\R^m$ through the quotient map
$q:\R^m\to\T^m$.
\end{lemma}

\begin{proof}
Let $F_n\subseteq(\T^d)^n$ be the fat wedge, the union of the sets on which
at least one coordinate is the identity.  Since $c|_{F_n}=0$, the map $c$
factors as a based continuous map
\[
\bar c:(\T^d)^n/F_n=(\T^d)^{\wedge n}\longrightarrow\T^m.
\]
Give $\T^d$ its standard based finite CW structure.  Every non-basepoint cell
has dimension at least one, so every non-basepoint cell of the $n$-fold smash
has dimension at least $n$.  For $n\ge2$ the smash product is connected and
has no one-cells; hence it is simply connected.  The universal covering
$q:\R^m\to\T^m$ gives a unique based lift
$\bar C:(\T^d)^{\wedge n}\to\R^m$.  Pulling back gives $C$.  Equivalently,
the lift restricts on the connected fat wedge to a lattice-valued map taking
the value $0$ at the basepoint, so it is normalized.
\end{proof}

\begin{lemma}[Haar contraction]\label{lem:haar}
Let $G$ be compact and let $V$ be a finite-dimensional real vector space with
trivial $G$-action.  Every continuous $V$-valued cocycle in positive degree is
a continuous coboundary.
\end{lemma}

\begin{proof}
Let $C\in C^n_{\cts}(G,V)$ satisfy $dC=0$, with $n\ge1$, and let $\mu$ be
normalized Haar measure.  Define
\[
B(g_1,\ldots,g_{n-1})
=(-1)^n\int_G C(g_1,\ldots,g_{n-1},t)\,d\mu(t).
\]
The function $B$ is continuous.  Integrate the identity
$dC(g_1,\ldots,g_n,t)=0$ in $t$.  Haar invariance changes the term containing
$g_nt$ into the same integral defining $B(g_1,\ldots,g_{n-1})$; collecting
the remaining terms gives $dB=C$.
\end{proof}

\section{Proof of the theorem}

\begin{proof}[Proof of Theorem~\ref{thm:main}]
In degree zero, the action is trivial, so the invariant group is $\T^m$.
In degree one, the cocycle identity says precisely that
$c(x+y)=c(x)+c(y)$.  Thus the cocycles are continuous homomorphisms
$\T^d\to\T^m$.  Degree-one boundaries vanish for a trivial action, and the
continuous homomorphisms between these tori are exactly the integer matrices
$\operatorname{Mat}_{m\times d}(\Z)$.

Now let $n\ge2$ and let $c:(\T^d)^n\to\T^m$ be a continuous cocycle.  By
Lemma~\ref{lem:normalization}, we may assume $c$ is normalized.  By
Lemma~\ref{lem:lift}, choose a normalized continuous lift
$C:(\T^d)^n\to\R^m$ with $q\circ C=c$.

Since $dc=0$, the map
\[
dC:(\T^d)^{n+1}\longrightarrow\R^m
\]
takes values in the deck lattice $\Z^m$.  Its domain is connected, so $dC$ is
constant.  At the all-identity tuple every term in $dC$ is a value of the
normalized cochain $C$ on a tuple containing an identity; all those terms are
zero.  Hence $dC=0$ everywhere.

Lemma~\ref{lem:haar} now gives a continuous cochain
$B:(\T^d)^{n-1}\to\R^m$ with $dB=C$.  Therefore
\[
c=q\circ C=d(q\circ B),
\]
so $c$ is a continuous coboundary.  This proves the asserted vanishing for
all $n\ge2$.
\end{proof}

\begin{remark}[Module action]
There is no omitted nontrivial toral action in Theorem~\ref{thm:main}.  The
group of continuous automorphisms of $\T^m$ is
$\operatorname{GL}(m,\Z)$ with the discrete topology.  Any continuous action
by automorphisms corresponds to a continuous homomorphism
$\T^d\to\operatorname{GL}(m,\Z)$, which is trivial because $\T^d$ is
connected.
\end{remark}

\begin{remark}[Why this does not misuse a long exact sequence]
The source emphasizes that the presentation
$\Z^m\to\R^m\to\T^m$ need not yield a long exact sequence for globally
continuous cochains, because a general continuous torus-valued cochain need
not admit a global continuous lift.  The proof above makes no such assertion.
It first normalizes a cochain of degree at least two; only then does the
fat-wedge factorization prove that this particular cochain has a lift.  The
cocycle equation is then checked again upstairs.  Thus the precise defect
highlighted in \cite{AustinPDE,AustinMoore} is bypassed, not assumed away.
\end{remark}

\section{Independent winding-number audit}

There is a direct degree-three check.  A map $\T^n\to\T$ has a winding vector
in $\Z^n$.  The bar differential induces
\[
D_2(b_1,b_2)=(-b_1,0,b_2),\qquad
D_3(a_1,a_2,a_3)=(0,a_2,a_2,0).
\]
Consequently
\[
\ker D_3=\{(a,0,c):a,c\in\Z\}=\operatorname{im}D_2.
\]
Thus every degree-three cocycle has coboundary winding, after which it lifts
to $\R$ and the proof above applies.  The accompanying exact-arithmetic script
checks these matrices, the bar signs, and $D_{n+1}D_n=0$ through degree eight.
This calculation is only a sanity check; the normalized smash-product proof
already handles all degrees and all finite-dimensional tori.

\section{Scope, novelty, and review status}

Theorem~\ref{thm:main} fully answers the exact Question 10.1 and strengthens
it to all degrees and all finite-dimensional source and coefficient tori.  It
does not complete Austin's broader program for continuous partial-difference
equations.  Nor does it assert vanishing for measurable, locally continuous,
smooth-local, Segal--Mitchison, or classifying-space cohomology; those are
different theories and may be nonzero \cite{AustinMoore,WagemannWockel}.

A bounded literature search on 13 August 2026 used the exact question and
notation, close variants with ``globally continuous cohomology'', circle and
torus coefficients, the source title/id/author, and citation searches.  The
primary sources \cite{AustinMoore,WagemannWockel} discuss the same globally
continuous theory but do not state this computation.  No later explicit
answer or equivalent full characterization was found.  Novelty confidence is
moderate: the argument may be implicit folklore under different notation,
but the exact published question and bounded search revealed no answer.

The result is marked \emph{candidate full solution, likely valid}.  Human
review should focus on Lemmas~\ref{lem:normalization} and~\ref{lem:lift}; once
those standard facts are accepted, the remaining proof is an explicit
cocycle and Haar-integration calculation.

\begin{thebibliography}{9}

\bibitem{AustinPDE}
T.~Austin,
\emph{Partial difference equations over compact Abelian groups, I: modules of
solutions}, arXiv:1305.7269, 2013; Question 10.1 on physical PDF page 91.

\bibitem{AustinMoore}
T.~Austin and C.~C.~Moore,
\emph{Continuity properties of measurable group cohomology},
arXiv:1004.4937, 2010.

\bibitem{WagemannWockel}
F.~Wagemann and C.~Wockel,
\emph{A cocycle model for topological and Lie group cohomology},
arXiv:1110.3304, 2011.

\end{thebibliography}

\end{document}
